
\section{No Difference}

Editorial article in “Folkebladet” for September 21, 1904.

In the last number of “Folkebladet” we discussed the remarkable claim advanced by the organ of the United Church, “Lutheraneren,” namely, that there is really no difference between the United Church and the Free Church, but that only “the leaders at Augsburg” have, in their cunning, “found out” something of that sort.

The question especially is this: Is the question, Do you want congregation? one such “dug-up” question, by which the “leaders” have wished to make it appear as though there were a difference where there was no difference?

All old Conference people, whether they belonged to the “old” or the “new” direction in the Conference, know well that there was a difference even within the Conference; and then they know all the more that there was a great difference between the “new” direction and the Synod.

But it may be that the anti-Missourians know nothing of this, or that they think that, in the days of conflict and union, they took such a decided stand with “the new direction,” that is, the Augsburgers within the Conference, that at any rate from that time there could be no talk of difference.

Believe that, whoever can!

We shall begin by presenting a testimony which shows that it has already become part of church history on a large scale that there was a difference in church principles and views, which lay at the foundation of the transfer question and of the coming-into-being of the Free Church.

We take this testimony from the most important theological work of the present age, namely Herzog’s Theological Realencyclopedia, issued in its third edition by Professor Hauck in Leipzig. In this work some 250 of the most learned men of our time within the Protestant churches are collaborating, and little nonsense is likely to find its way into it. In the fourteenth volume of this encyclopedia (printed in 1904) there is an article on the Lutheran Church in the United States, written by Adolf Spaeth, in which it says concerning the Lutheran Free Church:

“At the newly effected formation of the United Norwegian Lutheran Church in America, Augsburg Seminary in Minneapolis, which had previously belonged to the Norwegian-Danish Conference, was to pass over to the newly founded union. But difficulties soon arose in regard to this, since the views maintained by Augsburg Seminary concerning the requirements for theological formation, the pastoral office, congregation, and church did not meet with the approval of the United Norwegian Lutheran Church. These differences led to a complete rupture. In the year 1893 the United Church removed its theological institution from Augsburg Seminary and founded a new theological school under Professor M. O. Bøckman. The following year the leaders of Augsburg Seminary, Sverdrup and Oftedal, then organized the Lutheran Free Church upon a strongly congregationalistic basis. According to its latest church almanac, it numbers 112 pastors, 300 congregations, 38,000 communicants.”

We have translated this whole passage, that is, all that is found about the Free Church, because we regard it as an altogether impartial testimony, and because it is the understanding of the matter which, for a long time at any rate, will count as the sure and reliable one in church history.

We will not say that we are absolutely satisfied with everything in Spaeth’s testimony; but we accept it as, in all essentials, correct with respect to the difference between the majority of the United Church and the Augsburgers, and with respect to the difference which “led to a complete rupture.” We also wish to add that we feel assured that Spaeth’s information was not gathered from any man in the Lutheran Free Church. Should we be mistaken in this, it still does not alter our opinion of the significance of the testimony for the present and the future. It will remain standing as church-historical truth through long ages, however much “Lutheraneren” may try to overturn it.

And why should it not? It does the United Church no injustice that it is said that it did not approve Augsburg’s principles, unless “Lutheraneren” means that these principles are the only right ones. But surely “Lutheraneren” does not mean that?

The matter is simply this, that Augsburg has always, at any rate since the days of the “Open Declaration,”\footnote{See Professor Oftedal’s decisive testimony in court regarding the “Open Declaration,” FIXME page 198.} maintained the same views concerning theological training, the pastoral ministry, congregation, and church which it maintains to this very day. 

Georg Sverdrup

